\pdfbookmark{Kurzfassung / Abstract}{Abstract}
\section*{Kurzfassung}
\begin{otherlanguage}{ngerman}
\ac{Li-Fi} oder \ac{VLC} entwickelt sich zu einer vielversprechenden Lösung für das Problem der Frequenzknappheit und bietet Vorteile wie hohe Datenraten, leidet jedoch unter der begrenzten Batteriekapazität. Diese Arbeit befasst sich eingehend mit den Grundlagen und verwandten Arbeiten vieler \ac{VLC}-Architekturen, wobei \ac{EH} von Signal-light berücksichtigt wird. Diese Arbeit befasst sich mit den Herausforderungen bei der Entwicklung energieeffizienter und energiegewinnender \ac{VLC}-Systeme durch den Einsatz von \ac{SLIPT}-Techniken, insbesondere Power Splitting und Space Splitting. Als Instrument der Wahl für \ac{SLIPT}-Methoden wird ein Vergleich zwischen Solar- und \ac{PD}s angeboten.

Es wird ein umfassender theoretischer Ansatz zum Aufbau eines \ac{VLC}-Systems erläutert, zusammen mit der Auswahl der Komponenten, der Simulation und dem Bau eines Prototyps. Kritische Probleme, die während des Prozesses und der Designbewertung auftraten, werden erwähnt.

Der Empfänger decodierte Daten bei moderaten Entfernungen und Datenraten effektiv, hatte jedoch Schwierigkeiten, ausreichend Energie für einen vollständig autonomen Betrieb zu erzeugen. Die Solarzelle übertraf Photodioden in ihrer \ac{EH}-Fähigkeit, reichte jedoch für einen kontinuierlichen Betrieb nicht aus. Diese Ergebnisse unterstreichen die Notwendigkeit einer verbesserten Photodetektoreffizienz, verfeinerter Energiemanagementschaltungen und möglicherweise anderer hybrider Ansätze.

\end{otherlanguage}
\vfill
\section*{Abstract}
\begin{otherlanguage}{english}
\ac{Li-Fi} or \ac{VLC} is emerging as a promising solution to solve spectrum scarcity and offers advantages of high data rates, but suffers from finite battery capacity. This thesis thoroughly discusses Fundamentals and related works of many \ac{VLC} architectures, where \ac{EH} from Signal-light is considered. This thesis addresses the challenges of developing energy-efficient and energy-harvesting \ac{VLC} systems by adopting \ac{SLIPT} techniques, especially power and space splitting. Comparison between solar and \ac{PD}s is provided as an instrument of choice for \ac{SLIPT} methodologies. 

A comprehensive theoretical approach to building an \ac{VLC} system is explained, along with component selection, simulation, and prototype building. Critical issues faced during the process and design evaluation are mentioned.

The receiver effectively decoded data at moderate distances and data rates but struggled to generate sufficient harvested power for fully autonomous operation. The solar cell outperformed photodiodes in \ac{EH} capability, but was insufficient for continuous operation. These findings highlight the need for improved photodetector efficiency, refined power management circuits, and potentially other hybrid approaches.
\end{otherlanguage}
\vfill
